Turbo l'escargot joue à un jeu sur une grille avec $2024$ lignes et $2023$ colonnes. Il y a des monstres
cachés dans $2022$ cases. Initialement, Turbo ne connaît l'emplacement d'aucun monstre, mais il sait qu'il
y a exactement un monstre dans chaque ligne à part dans la première ligne et dans la dernière ligne, et que
chaque colonne contient au plus un monstre.

Turbo enchaîne plusieurs tentatives pour aller de la première ligne à la dernière ligne. À chaque tentative,
il choisit la case de la première ligne sur laquelle il commence, et se déplace de cases en cases, en allant
à chaque fois d'une case à une case ayant un côté en commun avec celle-ci. (Il a le droit de repasser sur une
case déjà visitée.) S'il accède à une case contenant un monstre, sa tentative se termine et il est transporté
sur la première ligne pour commencer une nouvelle tentative. Les monstres ne bougent pas, et Turbo se
souvient si chaque case qu'il a visitée contient un monstre ou non. S'il atteint n'importe quelle case de la
dernière ligne, sa tentative se termine et le jeu est fini.

Déterminer la valeur minimale de $n$ pour laquelle Turbo a une stratégie qui lui assure d'atteindre la
dernière ligne à sa $n$\textsuperscript{e} tentative ou avant, quelle que soit de la localisation des monstres.
